\documentclass[12pt]{article}
\usepackage{amsmath,amssymb,booktabs,longtable,array}
\usepackage{fontspec}
\usepackage{polyglossia}
\setmainlanguage{arabic}
\setotherlanguage{english}
\newfontfamily\arabicfont[Script=Arabic]{Amiri}
\usepackage{geometry}
\geometry{margin=2.1cm}
\usepackage{hyperref}

\title{بناء نموذج نقطي يعيد إنتاج منحنى الارتباط الكامل \(\rho(k)\)}
\author{مشروع تجربة البحث عن الأعداد الأولية}
\date{2026-09-01}

\begin{document}
\maketitle

\section{السؤال}

بعد المقارنة المباشرة مع نظرية الفترات القصيرة وLeung 2026، أصبح السؤال:

\begin{quote}
هل نستطيع تعديل النموذج النقطي بحيث يعيد إنتاج منحنى الارتباط الزوجي
\[
\rho(k),\qquad k=1,2,\ldots,
\]
قبل أن نضيف أي معلومة ثلاثية من Hardy--Littlewood \(k=3\)؟
\end{quote}

الهدف ليس ضبط كل lag على حدة. نريد نموذجًا زوجيًا فقط، قليل المعاملات،
ثم نسأل هل يظهر شكل منحنى الارتباط كاملًا.

\section{البيانات}

استعملنا نفس العينة:
\[
p_n\ge 4001,
\]
وعددها \(450\) فترة:
\[
p_n^2<x<p_{n+1}^2.
\]

ركزنا في النموذج النهائي على عجلة أساس:
\[
q_0=7,
\]
ثم random-residue sieve حتى:
\[
Q=7919.
\]

اختيار \(q_0=7\) لم يأت من ملاءمة منحنى \(\rho(k)\) في هذه التجربة؛
بل من المسح السابق، حيث كانت العجلات الصغيرة تترك للذيل العشوائي
أقوى إشارة زوجية.

\section{مرجع Leung: شكل \(\Delta(k)\)}

في حالة الفترات المتساوية، شكل التغاير الزوجي في نظرية Leung تحكمه الدالة:
\[
\Delta(t)
=
\frac12\left[
(t+1)\log(t+1)
-2t\log t
+(t-1)\log(t-1)
\right].
\]

ولمقارنة \emph{الشكل} فقط، نستخدم:
\[
D(k)
=
\frac{\Delta(k)}{\Delta(1)}.
\]

أما لأن فترات مربعات الأوليات ليست متساوية، فنستخدم أيضًا التنبؤ الأدق
المستخرج من قانون التباين:
\[
V_X(H)=H\log\left(\frac{X}{H}\right),
\]
ومن هوية التغاير بين فترتين منفصلتين:
\[
\mathrm{Cov}(I,J)
=
\frac{
V(d-a)+V(c-b)-V(c-a)-V(d-b)
}{2}.
\]

هذا يعطينا kernel نظريًا لكل زوج من فتراتنا، ثم نأخذ متوسطه عند كل lag.

\section{أين كان الخلل في النموذج القديم؟}

النموذج القديم كان:

\begin{enumerate}
\item عجلة ثابتة؛
\item random residue sieve في الذيل؛
\item ثم اختيار كل ناجٍ باستقلال باحتمال \(\theta_i\).
\end{enumerate}

إذا كان عدد الناجين بعد الغربلة العشوائية هو \(S_i\)، فإن الاختيار المستقل يعطي:
\[
K_i\mid S_i
\sim
\mathrm{Binomial}(S_i,\theta_i).
\]

ومن ثم:
\[
\mathrm{Cov}(K_i,K_j)
=
\theta_i\theta_j\mathrm{Cov}(S_i,S_j),
\qquad i\ne j,
\]
لكن:
\[
\mathrm{Var}(K_i)
=
\theta_i^2\mathrm{Var}(S_i)
+
\theta_i(1-\theta_i)\mathbb E[S_i].
\]

الحد الثاني في التباين هو ضوضاء مستقلة جديدة.
إذن نزيد التباين من غير أن نزيد التغاير بين الفترات، فتضعف:
\[
\rho(k).
\]

وهذا تفسير رياضي مباشر لما رأيناه سابقًا.

\section{تجربة block-quota thinning}

لنتأكد من ذلك بنموذج نقطي حقيقي، قسمنا ناجي كل فترة إلى كتل حجمها \(g\).
إذا كانت الكتلة تحتوي \(m\) ناجيًا، نحافظ على:
\[
\left\lfloor \theta_i m+U\right\rfloor,
\qquad U\sim\mathrm{Uniform}(0,1).
\]

عندما:
\[
g=1,
\]
نستعيد Bernoulli thinning المستقل.

وعندما يصبح \(g\) كبيرًا جدًا، نحصل تقريبًا على quota واحدة لكل الفترة،
أي نضيف قدرًا ضئيلًا جدًا من ضوضاء العد المستقلة.

المسح:
\[
g=1,2,4,8,\ldots,32768
\]
أظهر المقايضة بوضوح.

عند \(g=4\):
\[
R_{\rm MC}=0.3835
\]
مقابل:
\[
R_{\rm obs}=0.3803,
\]
أي أن التشتت أصبح ممتازًا.

لكن:
\[
\rho_1^{\rm MC}=-0.0864
\]
مقابل التنبؤ الزوجي:
\[
\rho_1^{\rm theory}=-0.1088.
\]

أما في حد quota الكامل تقريبًا:
\[
\rho_1^{\rm raw}\approx-0.1068,
\]
وهو قريب جدًا من النظرية، لكن:
\[
R_{\rm raw}\approx0.3085,
\]
أي أن التشتت أصبح أصغر من المطلوب.

إذن وجدنا قيدًا بنيويًا:

\begin{quote}
إذا استعملنا random-sieve signal الحالي ثم أضفنا ضوضاء عد مستقلة
لكي نرفع \(R\)، فإننا بالضرورة نضعف \(\rho(k)\).
\end{quote}

\section{هل random sieve نفسه يملك الشكل الصحيح؟}

نعم، وهذه أهم نتيجة وسيطة.

قبل إضافة ضوضاء thinning، أعطى random sieve عند \(q_0=7\):
\[
\rho_1\approx-0.1068,
\]
بينما نظرية الفترات القصيرة تعطي:
\[
-0.1088.
\]

وعند ملاءمة معامل واحد فقط في:
\[
\rho(k)\approx-a\Delta(k),
\]
كان:
\[
a_{\rm raw}\approx0.1495,
\]
بينما منحنى النظرية لفتراتنا يعطي:
\[
a_{\rm theory}\approx0.1587.
\]

أما خطأ الشكل بعد التطبيع عند lag 1 فكان:
\[
\mathrm{RMSE}_{\rm shape}\approx0.045.
\]

إذن random residue sieve يولد بالفعل kernel زوجيًا قريب الشكل والحجم
من kernel الفترات القصيرة.

المشكلة الأساسية كانت طريقة إضافة التباين بعد ذلك.

\section{النموذج الجديد: Pair-kernel quota model}

بما أن العجز هو تباين إضافي نحتاجه من دون إتلاف covariance،
أضفنا متغيرًا عشوائيًا زوجيًا واحدًا فقط.

نبدأ من quota count البنيوي:
\[
\theta_iS_i.
\]

ثم نضيف:
\[
\sqrt{\lambda V_i}\,G_i,
\]
حيث:
\[
G\sim N(0,R_{\rm pair}),
\]
و\(R_{\rm pair}\) هي مصفوفة الارتباط المحددة بالكامل من قانون الفترات القصيرة
حتى lag 50.

إذن:
\[
K_i^\ast
=
\theta_iS_i
+
\sqrt{\lambda V_i}\,G_i.
\]

ثم نقرب \(K_i^\ast\) عشوائيًا إلى عدد صحيح، ونختار بالضبط هذا العدد من
ناجي الغربال في الفترة.

المهم:

\begin{itemize}
\item لم نضع أي معامل ثلاثي.
\item لم نضبط \(\rho_1,\rho_2,\ldots\) كلًا على حدة.
\item الشكل الكامل لمصفوفة \(G\) مأخوذ من pair theory فقط.
\item المعامل الحر الوحيد \(\lambda\) حُدد من نقص التباين الهامشي، لا من الارتباط.
\end{itemize}

حسبنا:
\[
\lambda
=
R_{\rm obs}-R_{\rm struct}
=
0.07186.
\]

أي أن المعامل الوحيد يقول فقط:
``أضف مقدار التباين الناقص''.

\section{نتيجة التشتت}

بعد إضافة pair-kernel field:

\[
R_{\rm obs}=0.38033,
\]
و:
\[
R_{\rm model}=0.38097.
\]

في \(4000\) محاكاة:
\[
p\approx0.981.
\]

إذن التشتت أصبح مطابقًا تقريبًا، من دون التضحية بالـkernel الزوجي.

\section{نتيجة \(\rho(1)\)}

المرصود في السلسلة الواحدة:
\[
\rho_1^{\rm obs}=-0.12369.
\]

نظرية الفترات القصيرة:
\[
\rho_1^{\rm theory}=-0.10875.
\]

النموذج النقطي الجديد:
\[
\boxed{
\rho_1^{\rm model}=-0.10729.
}
\]

والقيمة المرصودة طبيعية داخل توزيع محاكاة النموذج:
\[
p\approx0.707.
\]

إذن النموذج لا يحاول نسخ الضوضاء الموجودة في قيمة lag واحدة؛
بل يطابق التنبؤ النظري الذي يمكن أن يولد تلك القيمة المرصودة.

\section{منحنى الارتباط الكامل}

قارنّا:
\[
k=1,\ldots,20.
\]

خطأ الشكل المطبع مقارنة بـ\(\Delta(k)\):
\[
\boxed{
\mathrm{RMSE}_{\rm shape}\approx0.037.
}
\]

ومعامل amplitude الملائم:
\[
a_{\rm model}\approx0.1515,
\]
مقابل:
\[
a_{\rm theory}\approx0.1587.
\]

إذن الشكل والحجم كلاهما قريبان.

أما اختبار كامل للمنحنى المرصود، باستعمال الانحرافات المعيارية للمحاكاة،
فأعطى:

\[
p_{\rm RMS}\approx0.667,
\]
و:
\[
p_{\rm max}\approx0.273.
\]

أي أن منحنى \(\rho(k)\) المرصود ككل ليس شاذًا تحت النموذج.

ظهرت قيم منفردة عند lags 18 و20 تبدو متطرفة إذا اختُبرت منفصلة،
لكنها لا تجعل المنحنى الكامل شاذًا بعد أخذ تعدد المسافات وتقلبها في الاعتبار.

\section{runs والثلاثيات: اختبار خارج الملاءمة}

لم نستخدم runs أو أي إحصائية ثلاثية في اختيار \(\lambda\).

مع ذلك:

\[
\mathrm{runs}_{\rm obs}=256,
\]
بينما:
\[
\mathrm{runs}_{\rm model}=241.04\pm10.13,
\]
مع:
\[
p\approx0.156.
\]

والثلاثيات الرتيبة:
\[
145
\]
مرصودة، مقابل:
\[
143.78\pm8.42
\]
في النموذج:
\[
p\approx0.906.
\]

كما أن جميع أنماط الإشارة الثلاثية:
\[
000,\ldots,111
\]
كانت غير شاذة، وأصغر قيمة \(p\) كانت نحو:
\[
0.227.
\]

إذن إضافة pair kernel الكامل أصلحت المقاييس التي كانت تدفعنا سابقًا
نحو التفكير في triples، من غير إدخال أي \(k=3\).

\section{ماذا أثبتت التجربة وماذا لم تثبت؟}

أثبتت التجربة على مستوى النموذج أن:

\[
\boxed{
\text{المعلومات الزوجية تكفي حاليًا لوصف جميع المقاييس التي اختبرناها.}
}
\]

لكن يجب الحذر.

الـpair-kernel correction أُدخل من نظرية الفترات القصيرة على مستوى
الفترات، ثم حُول إلى quota نقطية.
إذن هذا ليس اشتقاقًا كاملًا من random residue sieve وحده.

هو \emph{bridge model}:

\[
\text{random sieve at point level}
+
\text{known pair covariance}
\longrightarrow
\text{point set with correct interval kernel}.
\]

لذلك النجاح لا يعني أننا انتهينا من بناء نموذج أوليات طبيعي من المبادئ
المحلية فقط.

بل حدد بدقة ما ينقص random sieve الحالي:
\[
\boxed{
\text{مصدر زوجي إضافي للتباين له نفس kernel الفترات القصيرة تقريبًا.}
}
\]

\section{القرار البحثي}

لا ننتقل إلى triples الآن.

السؤال التالي أصبح أكثر دقة:

\begin{quote}
هل نستطيع توليد مقدار
\[
\lambda\approx0.0719
\]
وpair covariance المقابل من آلية نقطية محلية نفسها، بدل إدخاله كحقل Gaussian
على مستوى الفترات؟
\end{quote}

إذا أمكن ذلك، نكون قد نقلنا نجاح Leung/short-interval covariance من
طبقة ``تصحيح نظري'' إلى آلية مولدة داخل النموذج النقطي.

أما إذا فشلت كل آليات الأزواج المحلية بعد ذلك، وبقيت إحصاءات ثلاثية شاذة،
فعندها يصبح \(k=3\) خطوة مبررة.

\section{ملفات التجربة}

\begin{itemize}
\item \texttt{point\_model\_full\_rho\_kernel\_stage\_summary.csv}
\item \texttt{block\_quota\_thinning\_sweep\_q7.csv}
\item \texttt{pair\_kernel\_point\_model\_summary\_q7.csv}
\item \texttt{pair\_kernel\_point\_model\_rho\_q7.csv}
\item \texttt{pair\_kernel\_point\_model\_global\_rho\_test\_q7.csv}
\item \texttt{pair\_kernel\_point\_model\_sign\_triples\_q7.csv}
\item \texttt{plots/point\_model\_full\_rho\_kernel\_comparison.svg}
\item \texttt{plots/point\_model\_delta\_shape\_comparison.svg}
\item \texttt{plots/point\_model\_variance\_correlation\_tradeoff.svg}
\end{itemize}

\end{document}
